% Figure 2: the value of targeting decreases in the budget.
% Computed exactly from Proposition 1 (no simulation): population of n = 75 with
% K, R0 ~ Pareto(alpha = 2, min 1), rho = 0, A = 1, seed 659 (numpy default_rng);
% value of targeting = (f_opt - f_unif) / (f_unif - f_0), in percent;
% b = budget / total baseline resources. Data: fig2-curve.dat, generated by
% verify_gap_rule.py's companion snippet (see state.md, 2026-09-05).
\begin{figure}[t]
\centering
\begin{tikzpicture}
\begin{axis}[
  width=0.82\textwidth, height=5.4cm,
  axis lines*=left,
  xmin=0, xmax=2, ymin=0, ymax=175,
  xtick={0, 0.5, 1, 1.5, 2},
  ytick={0, 50, 100, 150},
  xlabel={budget scale $b$},
  ylabel={value of targeting (\%)},
  ylabel style={font=\small}, xlabel style={font=\small},
  tick label style={font=\small},
  axis line style={line width=0.4pt},
  tick style={line width=0.4pt, black},
  clip=false,
]
\addplot[black, line width=0.9pt, smooth] table {fig2-curve.dat};
\end{axis}
\end{tikzpicture}
\caption{The value of targeting decreases in the budget. The value of targeting is
the additional expected output from the optimal allocation over uniform funding, as
a fraction of uniform funding's gain over no funding; the budget scale $b$ is the
budget as a fraction of the population's total baseline resources. The curve shows
one population of $n = 75$ researchers drawn from the model's default distributions
($\alpha_K = \alpha_R = 2$, $\rho = 0$, $A = 1$); the value of targeting decreases
toward zero as the budget increases (Corollary~\ref{cor:targeting-vanishes}).}
\label{fig:targeting-value}
\end{figure}
